\documentclass[11pt,a4paper]{article}
\usepackage[margin=24mm]{geometry}
\usepackage{amsmath,amssymb}
\usepackage{booktabs}
\usepackage[T1]{fontenc}
\usepackage{lmodern}
\usepackage{microtype}
\usepackage{xcolor}
\usepackage{listings}
\usepackage{tikz}
\usepackage[hidelinks]{hyperref}
\usetikzlibrary{positioning,arrows.meta,fit,backgrounds}

\definecolor{ink}{HTML}{1A1A1A}
\definecolor{grey}{HTML}{5F5B54}
\definecolor{pale}{HTML}{F5F3EE}
\definecolor{down}{HTML}{1F4E8C}
\definecolor{up}{HTML}{1F6F4A}
\definecolor{warn}{HTML}{9C3B1B}

\lstdefinestyle{ts}{basicstyle=\ttfamily\footnotesize, showstringspaces=false,
  columns=fullflexible, keepspaces=true, breaklines=true,
  commentstyle=\color{grey}\itshape,
  backgroundcolor=\color{pale}, frame=none, xleftmargin=4mm, framexleftmargin=4mm}
\lstset{style=ts}
\newcommand{\note}[1]{\medskip\noindent\textbf{\textcolor{warn}{#1}}\quad}

\title{\vspace{-2em}\textbf{How we faked it}\\[2mm]
  \large Async event handlers as container morphisms:\\
  composing shapes forwards while responses flow back}
\date{\normalsize\today}

\begin{document}
\maketitle
\thispagestyle{empty}

\begin{abstract}
\noindent
The fake is not a type-level trick with a network attached. \textbf{The fake is
the async event handler.} A container's inhabitant $(p : P) \to R\,p$ is a
handler woken by a prompt; a morphism between containers sends shapes
\emph{forward} and brings responses \emph{back}; and \texttt{await} is what
allows both directions to be written as one straight-line function, in a server
that is simultaneously a client. This note is about that, and about what the
type system does and does not contribute to it.
\end{abstract}

\section{A container is an event handler}

\S3 of Ghani's note observes that event-driven programming ``has a feature that
turns out to be exactly a dependent type'': which responses are appropriate
depends on which event arrived. A container $(P, R)$ is that feature written
down --- $P$ the prompts, $R\,p$ the positions, the responses appropriate to $p$
--- and an inhabitant of the interface is the dependent function
$(p : P) \to R\,p$.

In this project that inhabitant is not a metaphor. It is the handler:

\begin{lstlisting}
serveContainer({
  name: "agri", port: 8802,
  parse:  (u: unknown) => AgPrompt | null,      // the prompt arrived as bytes
  answer: (p: AgPrompt, depth) => Promise<...>, // <-- THIS is  (p : P) -> R p
});
\end{lstlisting}

\noindent \texttt{answer} is woken by a socket, receives one prompt, and owes
exactly the response type that prompt indexes. Five such handlers are running
while the game is played.

\section{Two directions}

A morphism of containers $(P,R) \to (Q,T)$ is a pair:
\[
  u : P \to Q
  \qquad\qquad
  f : (p : P) \to T\,(u\,p) \to R\,p
\]

\noindent $u$ turns my prompt into one for the level below --- \textbf{shapes
travel forwards}. $f$ turns their response into mine --- \textbf{responses travel
back}. Control descends; data ascends. In a game about a command economy that
is not a coincidence: Gosplan does not reap anything. It converts ``feed the
towns'' into ``Narkomzem: harvest'', and converts the reply back into a number
for its own ledger.

The tower prints both directions as it runs. One \texttt{Sow}, across four
processes:

\begin{lstlisting}
:8801  |- gosplan   Sow order={"procurement":"firm","labour":{...
:8801  |   |- gosplan   ((agri (x) smelter) (x) transport)
:8802  |     |- agri      Harvest workforce={"fields":97,...} tractors=0 year=1928
:8802  |     -| agri      {grain, grade, perWorker, withoutTractors}
:8803  |     |- industry  Smelt workers=16 millCapacity=8 year=1928
:8803  |     -| industry  {steel, idleCapacity}
:8804  |     |- transport Lay workers=7 steel=10 year=1928
:8804  |     -| transport {added, steelUsed}
:8801  -| gosplan   {kind, year, harvest, steel, ...}
\end{lstlisting}

\noindent \verb#|-# is a shape going forward, \verb#-|# a response coming back.
Every indent is a morphism.

\section{\texttt{await} composes the two directions}

Here is the whole of a morphism, in code:

\begin{lstlisting}
const reply = await lower.run(u(prompt));   // shapes forward
return f(prompt, reply);                    // responses back
\end{lstlisting}

\noindent Both directions, one function body, read top to bottom.

\texttt{await} is what makes that possible. A handler is woken by a prompt, and
before it can answer it usually needs answers from below. At the \texttt{await}
it \emph{suspends}: it hands control back to the runtime, which is then free to
wake other handlers, and it resumes at that same point with the reply in hand.
So $u$ is applied on the line going down, $f$ on the line coming back, and the
suspension between them is invisible in the text of the function.

\note{Why this closes the algebra.} A morphism written this way is an ordinary
\texttt{async} function from a prompt to its response --- which is precisely what
a container's inhabitant is. So composing two morphisms is just calling one
inside the other, and the composite is again an \texttt{async} function from a
prompt to its response. \textbf{The morphisms are closed under composition in
the code, in the same way they are closed under composition in the theory}, and
a five-level tower is written as five ordinary functions calling one another.

\note{And it is why a server can also be a client.} Gosplan is woken by a socket
and, while answering, delegates to four commissariats. Because each of those
delegations suspends rather than occupies the runtime, Gosplan remains a
listening server throughout --- one handler, mid-answer, with four prompts
outstanding beneath it.

\section{The combinators, in both directions}

Each combinator is a way of composing the forward flow, and each has a matching
rule for the backward one.

\begin{center}
\footnotesize
\begin{tabular}{@{}cll@{}}
\toprule
& \textbf{shapes forward} & \textbf{responses back} \\
\midrule
$\boxtimes$ & a prompt to each side, dispatched together & a reply from \emph{each}; the position is a pair \\
$\times$    & a prompt to each side                      & a reply from \emph{exactly one} \\
$+$         & a prompt to one side                       & that side's position, unchanged \\
$\lhd$      & the second prompt is \emph{built from the first reply} & a pair of replies \\
\bottomrule
\end{tabular}
\end{center}

\subsection*{$\boxtimes$ --- forward together}

\begin{lstlisting}
const [left, right] = await Promise.all([a.run(s.left, d), b.run(s.right, d)]);
return { left, right };
\end{lstlisting}

\noindent Both shapes existed before the call, so neither can depend on the
other's reply --- independence is enforced by the shape type
$\{\texttt{left}, \texttt{right}\}$ itself: both prompts had to exist before the
call. The \texttt{Promise.all} sends them together, so the tensor costs one
round trip rather than two.

\subsection*{$\lhd$ --- where the two directions interleave}

This is the one that needs the hinge, and the reason the sequential composition
is the hard combinator:

\begin{lstlisting}
const first  = await a.run(s.first, depth);          // forward, then BACK
const second = await b.run(s.next(first), depth);    // the reply BECOMES the
                                                     // next shape: forward again
\end{lstlisting}

\noindent \texttt{s.next(first)} is the whole content of $\lhd$. A response
flowing back turns around and becomes a shape flowing forward. The composite's
shape is therefore not a pair of prompts but a prompt \emph{together with a
function} from its replies to the next prompt:

\begin{lstlisting}
step<SA extends Shape<A>, SB extends Shape<B>>(
  first: SA,
  next: (r: PosAt<A, SA>) => SB,          // SA inferred from `first`, so `next`
)                                          // is typed at THAT fibre alone
\end{lstlisting}

\noindent In the game: what you may \emph{sell} depends on what the railway
actually \emph{moved}, so the trade delegation is asked its question only once
the haul has answered.

\section{The tower makes the chain real}

$u$ is not a function call. It is a CLI tool:

\begin{lstlisting}
delegate.ts <port> <json-prompt> <depth>      # stdout is the reply
\end{lstlisting}

\noindent It knows nothing about any container --- port, JSON, depth, and it
carries the text across. That ignorance is the point: $u$ is a map of shapes,
and shapes are what survive a wire. A CLI tool calls a CLI tool, whose server
calls another, and the responses compose back up the same chain; the
\texttt{depth} header is what indents the trace above.

\note{Why processes rather than function calls.} Each commissariat holds a
hidden efficiency and a bias --- how much it inflates a shortfall --- and the
player plans against the reports, not the truth. Inside one process that privacy
is a convention, one stray import from collapsing. Across five processes with
only JSON between them the operating system enforces it: there is no reference
to reach through. The game's subject is an information asymmetry, and the wire
is what makes it a fact rather than a promise, at about a fifth of a second per
hop.

\section{What the type system contributes}

The handler is the fake; the types are what hold it to its word.

\subsection*{The container as its graph}

A container is presented as the union of its fibres, which over a
\emph{finite} index is exactly its $\Sigma$:
\[
  \Sigma\,(s : S).\,B\,s \;=\; \bigcup_{s \in S} \{s\} \times B\,s
\]

\begin{lstlisting}
interface Fib<S, P> { readonly shape: S; readonly pos: P }
type PosAt<C extends Cont, S> =
  C extends Fib<infer S1, infer P> ? (S extends S1 ? P : never) : never;
\end{lstlisting}

\noindent \texttt{Fib} is never constructed; it is taken apart by \texttt{infer}.
The handler's signature then reads \texttt{run<S extends Shape<C>>(shape: S):
Promise<PosAt<C, S>>} --- ask for a harvest, be handed a harvest report, uncast.

\subsection*{The finite index is what makes it exact}

That equality is the enabling condition, and it is worth being precise about
why it holds. The general definition of a container, as it appears in Agda,
\[
  \textbf{record } \mathrm{Container} : \quad
  \mathrm{Shape} : \mathrm{Set}, \qquad
  \mathrm{Pos} : \mathrm{Shape} \to \mathrm{Set}
\]
asks for $\mathrm{Pos}$ to be a function from \emph{values} to \emph{types} ---
a genuine dependent function, which is why containers are a dependently typed
idea and why the canonical example is the list functor, $\mathrm{Shape} =
\mathbb{N}$ with $\mathrm{Pos}\,n = \mathrm{Fin}\,n$.

Every container in this project has \emph{finitely} many prompts, because a
shape type is a tagged union. Over a finite index the graph presentation is not
an approximation of $\Sigma$ but an equality, and a conditional type
distributing over a union computes it exactly. That is the whole reason the
encoding is sound: not that TypeScript is cleverer than it looks, but that the
containers being encoded are small enough to write down fibre by fibre.

The same condition governs the same construction elsewhere. In Haskell the
finite case is a GADT indexed by its response type --- \texttt{Reconcile ::
(Int,Int) -> Request ReconReport} --- and \texttt{answer :: Request r -> r} is
the dependent function, checked branch by branch with no casts at all.

\subsection*{The bill}

Ten casts, two per combinator, because every \texttt{run} body destructures a
shape it knows through a conditional type and re-asserts the result. All ten sit
inside \texttt{lib/algebra.ts}; none is in the game, where no call site ever
asserts the type of a position.

Two pieces of machinery earn their keep quietly. \texttt{PosAt} tests
\texttt{S extends S1} --- the fibre a held shape \emph{lies in}, since a shape in
hand is usually narrower than its fibre's declared type. And \texttt{Runner}
carries a phantom \texttt{readonly fibres?: C}, never assigned, which gives
\texttt{C} an inferable position so that \texttt{seqC(a, b)} recovers its
arguments.

\section*{Summary}

\begin{center}
\footnotesize
\begin{tabular}{@{}p{0.26\textwidth}p{0.66\textwidth}@{}}
\toprule
\textbf{The fake} & an \texttt{async} event handler: $(p : P) \to R\,p$, woken by a socket \\
\textbf{What it buys} & shapes compose \emph{forwards} through $u$ and responses
  compose \emph{backwards} through $f$, in one straight-line function \\
\textbf{Why it composes} & a morphism is an ordinary \texttt{async} function from
  prompt to response, so composing two is calling one inside the other \\
\textbf{Why it scales} & suspension rather than occupation, so a handler can be
  mid-answer with four prompts outstanding and still be listening \\
\textbf{The rich combinator} & $\lhd$, where a response turns around and becomes
  the next shape \\
\textbf{What types add} & the response type follows from the prompt, uncast, at
  ten casts confined to one library \\
\textbf{What makes it exact} & finitely many prompts per container, so the graph
  presentation of $\Sigma$ is an equality \\
\end{tabular}
\end{center}

\end{document}
